% !TEX program = xelatex
% ============================================================
%  اتفاقية عدم المنافسة — Non-Compete Agreement Template
%  نموذج اتفاقية عدم منافسة تحدّ من نشاط الطرف في مجال عمل الطرف الآخر
%  Compile with: xelatex main.tex (run twice)
% ============================================================
%  Features:
%    - أطراف ومقدمات (Parties & Recitals)
%    - مواد مرقمة: تعريف النشاط المنافس، المدة، النطاق الجغرافي، الالتزامات، الاستثناءات، العلاج القانوني
%    - جدول النطاق الجغرافي والمدة (Territory & Duration Table)
%    - كتلة توقيعات للأطراف (Signatures Block)
%    - نص عربي قانوني مع مصطلحات إنجليزية بين قوسين
% ============================================================

\documentclass[12pt,a4paper]{article}

\usepackage[a4paper,margin=2.5cm,top=2.5cm,bottom=2.5cm,headheight=16pt]{geometry}
\usepackage{graphicx}
\usepackage{array}
\usepackage{booktabs}
\usepackage{tabularx}
\usepackage{multirow}
\usepackage{enumitem}
\usepackage{float}
\usepackage{titlesec}
\usepackage{fancyhdr}
\usepackage{setspace}
\usepackage{xcolor}

\usepackage{bidi}
\usepackage[no-math]{fontspec}
\usepackage[quiet]{polyglossia}

\setmainfont[Scale=1]{Amiri-Regular.ttf}
\newfontfamily\arabicfont[Script=Arabic,Scale=0.95]{Amiri-Regular.ttf}
\newfontfamily\englishfont[Scale=0.95]{TeX Gyre Termes}

\setdefaultlanguage[calendar=gregorian,hijricorrection=1]{arabic}
\setotherlanguage[variant=british]{english}

\usepackage[breaklinks,hidelinks]{hyperref}
\hypersetup{pdftitle={اتفاقية عدم المنافسة},pdfauthor={اسم الشركة}}

\titleformat{\section}{\large\bfseries}{\thesection}{0.7em}{}
\titlespacing*{\section}{0pt}{1.2ex plus 0.3ex}{0.8ex plus 0.2ex}

\pagestyle{fancy}
\fancyhf{}
\fancyhead[R]{\small\itshape اتفاقية عدم المنافسة (Non-Compete Agreement)}
\fancyhead[L]{\small اسم الشركة}
\fancyfoot[C]{\small\thepage}
\renewcommand{\headrulewidth}{0.3pt}

\newcolumntype{L}[1]{>{\raggedright\arraybackslash}p{#1}}
\newcolumntype{C}[1]{>{\centering\arraybackslash}p{#1}}

\setstretch{1.3}

\begin{document}

\begin{center}
{\LARGE\bfseries اتفاقية عدم المنافسة}\\[0.5em]
{\large (\LR{Non-Compete Agreement})}\\[1em]
\end{center}

\vspace{0.5em}

% TODO: استبدل التاريخ والمرجع بالقيم الفعلية
\noindent\textbf{رقم الاتفاقية:} \rule{4cm}{0.4pt} \hfill \textbf{التاريخ:} \rule{4cm}{0.4pt}

\vspace{1em}

% ============================================================
% PARTIES
% ============================================================
\noindent\textbf{الأطراف (\LR{Parties}):}

\noindent وقّع الطرفان أدناه على هذه الاتفاقية:

\begin{enumerate}
    \item \textbf{الطرف الأول (الشركة — \LR{Company}):} \rule{6cm}{0.4pt}، المسجّلة في \rule{4cm}{0.4pt}، رقم السجل التجاري: \rule{3cm}{0.4pt}، ويمثلها في التوقيع \rule{5cm}{0.4pt}.

    \item \textbf{الطرف الثاني (المُلتزِم — \LR{Restricted Party}):} \rule{6cm}{0.4pt}، رقم الهوية/الجواز: \rule{3cm}{0.4pt}، العنوان: \rule{5cm}{0.4pt}.
\end{enumerate}

\vspace{1em}

% ============================================================
% RECITALS
% ============================================================
\noindent\textbf{الديباجة (\LR{Recitals}):}

\begin{enumerate}
    \item تمتلك الشركة أعمالاً ومصالح تجارية في مجال \rule{6cm}{0.4pt}.
    \item سيتمكّن الطرف الثاني من الاطلاع على معلومات سرية (\LR{Confidential Information}) وعلاقات تجارية للشركة.
    \item ترغب الشركة في حماية مصالحها التجارية من المنافسة (\LR{Competition}) من قبل الطرف الثاني.
    \item يوافق الطرف الثاني على الالتزامات الواردة في هذه الاتفاقية مقابل \rule{4cm}{0.4pt} (الاعتبار/التعويض).
\end{enumerate}

\vspace{1em}

% ============================================================
% 1. DEFINITION OF COMPETITIVE BUSINESS
% ============================================================
\section{المادة الأولى: تعريف النشاط المنافس (\LR{Definition of Competitive Business})}
\label{sec:definition}
يُقصد بـ«النشاط المنافس» (\LR{Competitive Business}) أي نشاط تجاري أو مهني يباشره الطرف الثاني ويكون مماثلاً أو منافساً لنشاط الشركة، ويشمل على سبيل المثال لا الحصر:
\begin{enumerate}
    \item تقديم منتجات أو خدمات مماثلة لمنتجات وخدمات الشركة.
    \item العمل لدى منافس مباشر للشركة في نفس المجال.
    \item تأسيس أو الاستثمار في كيان ينافس الشركة.
    \item مساعدة أي طرف ثالث في منافسة الشركة.
\end{enumerate}

% ============================================================
% 2. DURATION
% ============================================================
\section{المادة الثانية: المدة (\LR{Duration})}
\label{sec:duration}
\begin{enumerate}
    \item \textbf{أثناء سريان العلاقة:} يلتزم الطرف الثاني بعدم المنافسة طوال فترة ارتباطه بالشركة.
    \item \textbf{بعد انتهاء العلاقة:} يستمر الالتزام بعد انتهاء العلاقة لمدة \rule{2cm}{0.4pt} شهراً/سنة.
    \item \textbf{في حال الإخلال:} يُمدَّد الالتزام بمدة تعادل مدة الإخلال.
\end{enumerate}

% ============================================================
% 3. TERRITORY
% ============================================================
\section{المادة الثالثة: النطاق الجغرافي (\LR{Territory})}
\label{sec:territory}
ينطبق الالتزام بعدم المنافسة على النطاق الجغرافي التالي:

% TODO: املأ النطاق الجغرافي والمدة الفعلية
\begin{table}[H]
\centering
\renewcommand{\arraystretch}{1.5}
\begin{tabularx}{\textwidth}{L{5cm} X L{3cm}}
\toprule
\textbf{البند} & \textbf{التفاصيل} & \textbf{المدة} \\
\midrule
النطاق المحلي & \rule{6cm}{0.4pt} & \rule{2.5cm}{0.4pt} \\
النطاق الإقليمي & \rule{6cm}{0.4pt} & \rule{2.5cm}{0.4pt} \\
النطاق الدولي & \rule{6cm}{0.4pt} & \rule{2.5cm}{0.4pt} \\
\bottomrule
\end{tabularx}
\end{table}

% ============================================================
% 4. OBLIGATIONS
% ============================================================
\section{المادة الرابعة: التزامات الطرف المُلتزِم (\LR{Obligations})}
\label{sec:obligations}
يلتزم الطرف الثاني بما يلي:
\begin{enumerate}
    \item \textbf{عدم المنافسة (\LR{Non-Compete}):} عدم مباشرة أي نشاط منافس خلال فترة الالتزام.
    \item \textbf{عدم استمالة العملاء (\LR{Non-Solicitation of Customers}):} عدم التواصل مع عملاء الشركة لعرض خدمات بديلة.
    \item \textbf{عدم استمالة الموظفين (\LR{Non-Solicitation of Employees}):} عدم استمالة موظفي الشركة أو تشجيعهم على ترك العمل.
    \item \textbf{عدم استمالة الموردين (\LR{Non-Solicitation of Suppliers}):} عدم التدخل في علاقة الشركة مع مورديها.
    \item \textbf{عدم استخدام الأسرار التجارية (\LR{No Use of Trade Secrets}):} عدم استخدام معلومات الشركة السرية لصالح نشاط منافس.
\end{enumerate}

% ============================================================
% 5. EXCEPTIONS
% ============================================================
\section{المادة الخامسة: الاستثناءات (\LR{Exceptions})}
\label{sec:exceptions}
لا يشمل الالتزام بعدم المنافسة الحالات التالية:
\begin{enumerate}
    \item \textbf{الملكية السلبية (\LR{Passive Ownership}):} امتلاك أقل من \rule{2cm}{0.4pt}\% من أسهم شركة عامة دون مشاركة في الإدارة.
    \item \textbf{النشاط غير المرتبط:} العمل في مجال لا علاقة له بنشاط الشركة.
    \item \textbf{بعد انتهاء المدة:} أي نشاط بعد انتهاء فترة الالتزام المتفق عليها.
    \item \textbf{بموافقة كتابية:} أي نشاط يوافق عليه الطرف الأول كتابةً.
\end{enumerate}

% ============================================================
% 6. REMEDIES
% ============================================================
\section{المادة السادسة: العلاج القانوني (\LR{Remedies})}
\label{sec:remedies}
\begin{enumerate}
    \item \textbf{التعويض (\LR{Damages}):} يحق للشركة المطالبة بالتعويض عن جميع الأضرار الناتجة عن الإخلال.
    \item \textbf{الأوامر القضائية (\LR{Injunction}):} يحق للشركة الحصول على أوامر قضائية عاجلة لمنع استمرار الإخلال.
    \item \textbf{الغرامة الاتفاقية (\LR{Liquidated Damages}):} يدفع الطرف الثاني غرامة قدرها \rule{3cm}{0.4pt} عن كل إخلال، دون المساس بالحق في التعويض الإضافي.
    \item \textbf{المصاريف القانونية (\LR{Legal Costs}):} يتحمّل الطرف المُخلّ جميع المصاريف القانونية وأتعاب المحاماة.
\end{enumerate}

% ============================================================
% 7. GOVERNING LAW
% ============================================================
\section{المادة السابعة: القانون الحاكم وحل النزاعات (\LR{Governing Law \& Dispute Resolution})}
\label{sec:law}
\begin{enumerate}
    \item تُفسَّر هذه الاتفاقية وفقاً لقوانين \rule{5cm}{0.4pt}.
    \item تُحلّ النزاعات (\LR{Disputes}) الناشئة عن هذه الاتفاقية عن طريق التحكيم في \rule{4cm}{0.4pt}.
    \item تُعتبر أحكام المحكّمين نهائية وملزمة للطرفين.
\end{enumerate}

% ============================================================
% SIGNATURES
% ============================================================
\vspace{2em}
\section*{التوقيعات (\LR{Signatures})}

% TODO: املأ بيانات الموقّعين
\begin{table}[H]
\centering
\renewcommand{\arraystretch}{2.5}
\begin{tabularx}{\textwidth}{L{5cm} X X}
\toprule
& \textbf{الطرف الأول (الشركة)} & \textbf{الطرف الثاني (المُلتزِم)} \\
\midrule
\textbf{الاسم} & --- & --- \\
\textbf{الصفة} & --- & --- \\
\textbf{التوقيع} & \rule{4cm}{0.4pt} & \rule{4cm}{0.4pt} \\
\textbf{التاريخ} & \rule{4cm}{0.4pt} & \rule{4cm}{0.4pt} \\
\bottomrule
\end{tabularx}
\end{table}

\end{document}
